% TEMPLATE-MILC-v3.tex - plantilla maestra para todas las guias MILC v3
% Ver scripts/generadores/build-guia-milc.py para la lista completa de
% claves esperadas. El script valida que no queden placeholders sin
% reemplazar antes de escribir el archivo final.

\documentclass[11pt,letterpaper]{article}
\usepackage[margin=1.75cm]{geometry}
\usepackage[table]{xcolor}
\usepackage{fontspec}
\usepackage[spanish]{babel}
\usepackage{tikz}
\usepackage[most]{tcolorbox}
\usepackage{tabularx}
\usepackage{array}
\usepackage{fancyhdr}
\usepackage{titlesec}
\usepackage{ragged2e}
\usepackage{enumitem}
\usepackage{microtype}

\setmainfont{Helvetica Neue}
\setsansfont{Helvetica Neue}
\setlength{\parindent}{0pt}
\setlength{\parskip}{6pt}
\hyphenpenalty=200
\exhyphenpenalty=200
\pretolerance=400
\tolerance=2000
\setlength{\emergencystretch}{1em}
\renewcommand{\arraystretch}{1.25}
\setlist{leftmargin=5mm,itemsep=1mm,topsep=1mm}
\newcolumntype{Y}{>{\RaggedRight\arraybackslash}X}

\definecolor{milcMagenta}{HTML}{E6007E}
\definecolor{milcVino}{HTML}{5A0038}
\definecolor{milcTurquesa}{HTML}{0093A5}
\definecolor{milcVerde}{HTML}{58A923}
\definecolor{milcMostaza}{HTML}{E5B400}
\definecolor{milcOcre}{HTML}{B86600}
\definecolor{milcNegro}{HTML}{241816}
\definecolor{milcGris}{HTML}{F7F7F4}
\definecolor{milcLinea}{HTML}{E8E2DD}
\definecolor{milcGradePrimary}{HTML}{0066FF}
\definecolor{milcGradeDark}{HTML}{003D99}
\definecolor{milcGradeSoft}{HTML}{D6E8FF}
\definecolor{milcGradeAccent}{HTML}{CFE7FF}
\definecolor{milcGradeText}{HTML}{FFFFFF}
\color{milcNegro}

\pagestyle{fancy}
\fancyhf{}
\lhead{\small Tecnología e Informática · Grado 11}
\rhead{\small Guía 4 · MILC v3}
\cfoot{\small\thepage}
\renewcommand{\headrulewidth}{0pt}

\newtcolorbox{titlebox}[1]{
  enhanced,colback=#1,colframe=#1,arc=7mm,boxrule=0pt,
  left=7mm,right=7mm,top=3mm,bottom=3mm,
  before skip=8pt,after skip=8pt,
  fontupper=\sffamily\bfseries\Large\color{white}
}

\newtcolorbox{softbox}[3]{
  enhanced,colback=#2,colframe=#1,borderline west={4pt}{0pt}{#1},
  arc=2mm,boxrule=.4pt,left=4mm,right=4mm,top=3mm,bottom=3mm,
  before skip=6pt,after skip=8pt,title={#3},coltitle=white,
  colbacktitle=#1,fonttitle=\sffamily\bfseries,fontupper=\RaggedRight,
  attach boxed title to top left={xshift=3mm,yshift=-2mm},
  boxed title style={arc=2mm, boxrule=0pt}
}

\newtcolorbox{drawbox}[2]{
  enhanced,colback=white,colframe=#1,arc=4mm,boxrule=1pt,
  left=5mm,right=5mm,top=4mm,bottom=4mm,before skip=8pt,after skip=8pt,
  title={#2},coltitle=white,colbacktitle=#1,fonttitle=\sffamily\bfseries,
  fontupper=\RaggedRight,
  attach boxed title to top left={xshift=4mm,yshift=-2mm},
  boxed title style={arc=3mm, boxrule=0pt}
}

\newtcolorbox{infoband}[1]{
  enhanced,colback=#1!8,colframe=#1!50,arc=2mm,boxrule=.5pt,
  left=4mm,right=4mm,top=2mm,bottom=2mm,before skip=4pt,after skip=4pt,
  fontupper=\small\RaggedRight
}

% Puente narrativo entre secciones (contrato editorial Conéctate).
% Da continuidad: "Acabas de X. Ahora vas a Y porque Z."
% Visualmente: banda izquierda en color del grado, texto en itálica pequeña.
\newtcolorbox{puentebox}{
  enhanced,colback=milcGradeSoft,colframe=milcGradeDark,
  borderline west={3pt}{0pt}{milcGradeDark},
  arc=0mm,boxrule=0pt,left=5mm,right=4mm,top=2.5mm,bottom=2.5mm,
  before skip=4pt,after skip=8pt,
  fontupper=\itshape\small\color{milcNegro}\RaggedRight
}
\newcommand{\puente}[1]{%
  \begin{puentebox}%
    \textcolor{milcGradeDark}{\textbf{\textrightarrow}}\hspace{2mm}#1%
  \end{puentebox}%
}

\newcommand{\linead}[1]{\noindent\color{milcLinea}\rule{#1}{0.5pt}\color{milcNegro}}
\newcommand{\respuesta}[1]{\par\vspace{2mm}\foreach \n in {1,...,#1}{\noindent\color{milcLinea}\rule{\linewidth}{0.5pt}\par\vspace{5mm}}\color{milcNegro}}
\newcommand{\checkbox}{\textcolor{milcGradeDark}{\fbox{\phantom{x}}}\hspace{5pt}}

\newcommand{\phasecard}[4]{
\begin{tcolorbox}[enhanced,colback=#3,colframe=#2,arc=3mm,boxrule=.5pt,height=3.05cm,valign=center,halign=center,left=1.5mm,right=1.5mm,top=2mm,bottom=2mm]
{\sffamily\bfseries\fontsize{22}{24}\selectfont\textcolor{#2}{#1}}\par
{\sffamily\bfseries\small #4}
\end{tcolorbox}
}

\begin{document}
\thispagestyle{empty}
\begin{tikzpicture}[remember picture,overlay]
  \fill[white] (current page.south west) rectangle (current page.north east);
  \path[fill=milcGradePrimary,rounded corners=16mm]
    ([xshift=.55cm,yshift=-.55cm]current page.north west) rectangle
    ([xshift=-.55cm,yshift=3.65cm]current page.south east);

  \node[anchor=north west,text=milcGradeText] at ([xshift=2.0cm,yshift=-1.85cm]current page.north west)
    {\sffamily\bfseries\fontsize{14}{16}\selectfont GRADO};
  \node[anchor=north west,text=milcGradeText,opacity=.82] at ([xshift=2.0cm,yshift=-2.75cm]current page.north west)
    {\sffamily\bfseries\fontsize{9}{11}\selectfont SERIE GUÍAS · TECNOLOGÍA E INFORMÁTICA};

  \begin{scope}[shift={([xshift=-3.05cm,yshift=-2.55cm]current page.north east)}]
    \fill[white,opacity=.80] (-.62,-.52) rectangle (.62,.52);
    \draw[milcGradeText,line width=1pt] (-.62,-.52) rectangle (.62,.52);
    \fill[milcGradeDark] (-.42,-.38) rectangle (-.22,.06);
    \fill[milcGradeAccent] (-.10,-.38) rectangle (.10,.24);
    \fill[milcGradePrimary!60!black] (.22,-.38) rectangle (.42,.38);
  \end{scope}

  \node[anchor=west,text=milcGradeText] at ([xshift=1.9cm,yshift=-12.85cm]current page.north west)
    {\sffamily\bfseries\fontsize{124}{124}\selectfont 11};
  \draw[milcGradeText,line width=1.7pt]
    ([xshift=4.52cm,yshift=-11.68cm]current page.north west) circle (.13cm);

  \node[anchor=west,text=milcGradeText,text width=8.7cm,align=left] at ([xshift=10.35cm,yshift=-11.6cm]current page.north west)
    {
     {\sffamily\bfseries\fontsize{19}{22}\selectfont Plantillas web ---\\elegir y adaptar\\sin perderse uno mismo}\\[4mm]
     {\sffamily\bfseries\fontsize{23}{26}\selectfont Guía 4-11-TIC}\\[2mm]
     {\sffamily\fontsize{13}{16}\selectfont Período 1 · Presencia y marca digital}\\[5mm]
     {\sffamily\bfseries\fontsize{11}{13}\selectfont Institución Educativa Sor María Juliana}\\[1mm]
     {\sffamily\fontsize{10}{12}\selectfont Autor: Dr. Álvaro Cárdenas Orozco}
    };

  \path[fill=milcGradeSoft,rounded corners=7mm]
    ([xshift=1.35cm,yshift=.85cm]current page.south west) rectangle
    ([xshift=-1.35cm,yshift=4.25cm]current page.south east);
  \draw[milcGradeDark,line width=.65pt,rounded corners=7mm]
    ([xshift=1.35cm,yshift=.85cm]current page.south west) rectangle
    ([xshift=-1.35cm,yshift=4.25cm]current page.south east);
  \node[anchor=center,text=milcNegro,text width=17.0cm,align=center] at ([yshift=2.55cm]current page.south)
    {\sffamily\fontsize{10}{12}\selectfont Metodología MILC · Apertura ancestral · Pensamiento computacional · Triángulo Dussel-Estoicismo-Floridi\\
    \textcolor{milcGradeDark}{\bfseries Producto:} Sitio publicado en GitHub Pages (o equivalente) con plantilla adaptada};
\end{tikzpicture}
\mbox{}
\newpage

\section*{Apertura: Plantillas web --- elegir y adaptar sin perderse uno mismo}

\begin{softbox}{milcGradeDark}{milcGradeSoft}{Saber ancestral}
Los oficios heredados se aprenden con \textbf{moldes}: la receta de la abuela, el corte del sastre, el patrón del tejedor, las décimas del Pacífico. El maestro no enseña con palabras: enseña con un molde que el aprendiz debe usar primero igual, y después adaptar a su mano. Lo que cambia es la firma; lo que se conserva es la estructura.
\end{softbox}

\begin{softbox}{milcTurquesa}{milcGris}{Saber contemporáneo}
En la web moderna existen \textbf{plantillas} en GitHub Pages, Carrd, Notion sites, Webflow, Astro starters. Son moldes profesionales que aceleran tu trabajo --- siempre y cuando los adaptes con criterio en lugar de dejarlos genéricos.
\end{softbox}

\begin{softbox}{milcMostaza}{milcGris}{Pregunta-puente}
\textit{¿Qué se mantiene y qué cambia cuando tú haces la receta de tu abuela? ¿Qué la sigue convirtiendo en \emph{esa} receta y no en una copia genérica?}
\end{softbox}

\begin{softbox}{milcVerde}{milcGris}{Saber y saber hacer}
Al terminar podrás: (1) \textbf{analizar} 3 plantillas similares e identificar qué tienen en común; (2) \textbf{explicar} qué se conserva y qué cambia al adaptar un molde; (3) \textbf{aplicar} una plantilla a tu identidad de marca y publicarla; (4) \textbf{evaluar} si tu sitio adaptado sigue siendo ``tuyo''.
\end{softbox}

\small
\begin{tabularx}{\textwidth}{>{\bfseries}p{3.0cm}Y}
\rowcolor{milcGradePrimary!12}Autor & Dr. Álvaro Cárdenas Orozco\\
DBA & Comunicación profesional en entornos digitales (MEN, Lineamientos T\&I)\\
Referentes & Floridi (infoética) · Dussel (estética de la liberación) · Estoicismo\\
Producto final & Sitio publicado en GitHub Pages (o equivalente) con plantilla adaptada\\
\end{tabularx}
\normalsize

\puente{Antes de empezar, mira el viaje completo. Hoy publicas tu sitio en internet --- por primera vez ya no será un archivo en tu computador, será una URL que cualquiera puede abrir.}

\begin{titlebox}{milcGradeDark}
Ruta MILC: así vamos a aprender
\end{titlebox}
\begin{tabularx}{\textwidth}{>{\centering\arraybackslash}X>{\centering\arraybackslash}X>{\centering\arraybackslash}X>{\centering\arraybackslash}X}
\phasecard{1}{milcVerde}{milcGris}{Escuta\\[-1mm]\small Observo, escucho y pregunto.} &
\phasecard{2}{milcTurquesa}{milcGris}{Sistematización\\[-1mm]\small Pensamiento computacional.} &
\phasecard{3}{milcMagenta}{milcGris}{Praxis\\[-1mm]\small Construyo y aplico.} &
\phasecard{4}{milcVino}{milcGris}{Evaluación\\[-1mm]\small Reflexión liberadora.}
\end{tabularx}


\puente{Antes de elegir tu plantilla, vas a aprender a mirar otras. Compararlas te enseña la gramática común de las plantillas modernas --- y revela qué espacio hay para hacerla tuya.}

\begin{titlebox}{milcVerde}
Fase 1 · Escuta
\end{titlebox}

\begin{softbox}{milcVerde}{milcGris}{Escena para escuchar}
\textbf{Actividad 1 · ANALIZA --- Tres plantillas, qué tienen en común} (15 min · individual). Visita 3 sitios construidos sobre plantillas (puedes mirar catálogos de GitHub Pages templates, Carrd showcase, o Astro themes). Para cada uno: ¿qué partes son intercambiables (logo, paleta, tipografía, fotos)? ¿qué partes son estructura fija (orden de secciones, jerarquía)?
\end{softbox}

\checkbox Revisé 3 plantillas distintas en uno de los catálogos sugeridos.\\
\checkbox Distinguí qué partes son intercambiables y cuáles son fijas.\\
\checkbox Tomé nota del nombre y URL de la plantilla que más se acerca a mi marca.

\begin{infoband}{milcVerde}
\textbf{Lo que va al cuaderno (Actividad 1 · ANALIZA).} Título: «Actividad 1 --- Tres plantillas comparadas». \textbf{Formato:} tabla de 4 columnas (Plantilla | URL | Partes intercambiables | Partes fijas), 3 filas. Al final un párrafo de 3-4 renglones eligiendo una plantilla y justificando por qué. \textbf{Sabes que terminaste cuando} tu elección es coherente con tu paleta y tipografía de la Guía 2.
\end{infoband}

\puente{Ya viste tres plantillas. Ahora vamos a entender el \textbf{principio} detrás de adaptar un molde sin perderse: qué se queda fijo, qué se cambia, y qué decisión convierte una plantilla genérica en \emph{tu} sitio.}

\begin{titlebox}{milcTurquesa}
Fase 2 · Sistematización con pensamiento computacional
\end{titlebox}

\begin{softbox}{milcTurquesa}{milcGris}{Cuatro pilares aplicados al tema}
Adaptar una plantilla es un acto de criterio, no de pereza. Vamos a descomponer ese criterio en cuatro decisiones que diferencian un sitio único de uno genérico.
\end{softbox}

\small
\begin{tabularx}{\textwidth}{>{\bfseries\RaggedRight\arraybackslash}p{3.6cm}Y}
\rowcolor{milcTurquesa!90}\color{white}Pilar & \color{white}\bfseries Aplicación al tema\\
1. Descomposición & Una adaptación se descompone en 4 decisiones: (1) qué \textbf{conservar} de la plantilla (estructura, layout), (2) qué \textbf{cambiar} (colores, fuentes, fotos, copy), (3) qué \textbf{agregar} (tu marca, tu voz), (4) qué \textbf{eliminar} (secciones que no te sirven).\\
2. Reconocimiento de patrones & Toda buena adaptación sigue el patrón 70-20-10: \textasciitilde 70\% queda como vino la plantilla (estructura), \textasciitilde 20\% se personaliza (colores, fuentes, copy), \textasciitilde 10\% se agrega o quita (lo único de tu marca). Cambia el patrón y el resultado se siente o copiado o frankenstein.\\
3. Abstracción & Lo esencial de tu adaptación cabe en una pregunta: \emph{¿en qué se nota que este sitio es tuyo y no de otro que use la misma plantilla?}. Si no tienes 3 respuestas, todavía no lo has hecho tuyo.\\
4. Algoritmo conceptual & Pasos: (1) haz fork de la plantilla en tu GitHub; (2) clónala y ábrela en VS Code; (3) cambia colores y fuentes por los de tu sistema visual; (4) reemplaza fotos y copy genéricos por los tuyos; (5) publica con GitHub Pages.\\
\end{tabularx}
\normalsize

\begin{drawbox}{milcTurquesa}{Anatomía de una buena adaptación}
\textbf{Conservar:} estructura semántica de la plantilla. No reinventes la rueda. \\[1mm] \textbf{Cambiar:} paleta y tipografía a las tuyas (de la Guía 2). Copy a tu voz. \\[1mm] \textbf{Agregar:} un detalle único de tu marca (un símbolo, un microcopy especial, una sección propia). \\[1mm] \textbf{Eliminar:} cualquier sección genérica que la plantilla traiga y que tú no necesites. \emph{Menos es más.}
\end{drawbox}

\begin{softbox}{milcMostaza}{milcGris}{Errores comunes y su corrección}
(1) Dejar copy genérico (``Welcome to my site!''): muerte instantánea. (2) Cambiar TODO y romper la coherencia visual de la plantilla. (3) Conservar fotos placeholder o nombres de prueba (``John Doe''). (4) No probar el sitio en móvil: la mayoría de tus visitantes te leerán ahí.
\end{softbox}

\begin{infoband}{milcTurquesa}
\textbf{Lo que va al cuaderno (Actividad 2 · EXPLICA).} Título: «Actividad 2 --- Plan de adaptación de mi plantilla». \textbf{Formato:} 4 columnas (Conservar | Cambiar | Agregar | Eliminar), cada una con 3-5 elementos específicos de tu plantilla elegida. \textbf{Sabes que terminaste cuando} las 4 columnas tienen contenido concreto (no ``el color'' sino ``el azul \#0066FF del navbar por mi naranja \#FF6B35'').
\end{infoband}

\puente{Tienes el criterio. Hora de ejecutar: eliges una plantilla, la adaptas a tu identidad (usando tu paleta y tipografía de la guía anterior) y la publicas en GitHub Pages.}

\begin{titlebox}{milcMagenta}
Fase 3 · Praxis con pensamiento computacional
\end{titlebox}

\begin{softbox}{milcMagenta}{milcGris}{Construyendo el producto con los cuatro pilares}
\textbf{Actividad 3 · APLICA --- Publicar tu sitio adaptado} (40 min · individual). Hora de ejecutar. Vas a clonar la plantilla, adaptarla con tu paleta y tu copy, y publicarla en GitHub Pages. Al final tendrás una URL pública que cualquiera puede abrir desde el mundo.
\end{softbox}

\small
\begin{tabularx}{\textwidth}{>{\bfseries\RaggedRight\arraybackslash}p{3.6cm}Y}
\rowcolor{milcMagenta!90}\color{white}Pilar & \color{white}\bfseries Aplicación al producto\\
Descomposición & Tu publicación se descompone en pasos: fork del repo → cambiar archivos clave (paleta, copy) → commit → push → activar GitHub Pages → verificar URL pública.\\
Patrones & Sigue el patrón de despliegue: rama main, configurar Pages desde Settings → Pages, espera 1-2 min, abre la URL en pestaña incógnita.\\
Abstracción & Cada commit expresa una sola idea: ``colores a mi paleta'', ``copy del header'', ``foto personal''. Si haces commits gigantes, no podrás revertir cambios específicos.\\
Algoritmo de armado & Algoritmo: (1) fork de la plantilla a tu GitHub; (2) clone local; (3) abrir en VS Code; (4) cambiar paleta (CSS variables o archivo de tema); (5) reemplazar copy y fotos; (6) commit y push; (7) Settings → Pages → main branch → save; (8) esperar URL; (9) verificar en móvil y desktop.\\
\end{tabularx}
\normalsize

\begin{drawbox}{milcMagenta}{Checklist antes de declarar el sitio publicado}
\checkbox La URL pública carga sin errores 404. \\[1mm] \checkbox Cero ``Lorem ipsum'' o copy de placeholder. \\[1mm] \checkbox Los colores son los de tu paleta (Actividad 3, Guía 2). \\[1mm] \checkbox La tipografía es tu par (Actividad 3, Guía 2). \\[1mm] \checkbox Se ve bien en móvil (revisa con DevTools a 360 px). \\[1mm] \checkbox Tu nombre y propuesta de valor aparecen en el hero.
\end{drawbox}

\begin{softbox}{milcMagenta}{milcGris}{Plantilla de trabajo}
\textbf{Plantilla elegida.} \textit{Nombre + URL.} \\[1mm] \textbf{Mi sitio.} \textit{La URL pública en GitHub Pages.} \\[1mm] \textbf{Cambios aplicados.} \textit{4 columnas como en Actividad 2.} \\[1mm] \textbf{Captura.} \textit{Pega o dibuja una captura de tu home en el cuaderno.}
\end{softbox}

\begin{infoband}{milcMagenta}
\textbf{Lo que va al cuaderno (Actividad 3 · APLICA).} Título: «Actividad 3 --- Mi sitio publicado». \textbf{Formato:} (1) URL pública en grande arriba, (2) tabla de cambios aplicados, (3) captura impresa o dibujada del home. \textbf{Sabes que terminaste cuando} la URL carga en móvil de un compañero, sin lorem ipsum, con tu paleta y tu tipografía.
\end{infoband}

\puente{Lo que sigue es el resumen de \emph{qué} entregas hoy: una URL pública con tu sitio adaptado. Es el primer ``hola mundo'' real de tu marca en internet.}

\begin{titlebox}{milcMostaza}
Producto de la sesión
\end{titlebox}
\textbf{Sitio personal publicado en GitHub Pages con plantilla adaptada}.
\begin{softbox}{milcMostaza}{milcGris}{Criterios de calidad}
(1) La URL pública es accesible para cualquiera con internet. (2) Cero contenido placeholder (texto, imágenes, nombres de prueba). (3) La paleta y la tipografía coinciden con tu sistema visual (Guía 2). (4) En móvil se lee y navega sin problemas. (5) Hay al menos UN detalle único (símbolo, sección propia, microcopy) que no estaba en la plantilla original.
\end{softbox}

\puente{Antes de cerrar, mira tu sitio publicado desde las cinco dimensiones humanas. Publicar tiene costo emocional y ciudadano --- estás haciendo pública tu identidad.}

\begin{titlebox}{milcVino}
Fase 4 · Evaluación liberadora — cinco dimensiones
\end{titlebox}

\begin{softbox}{milcVino}{milcGris}{Sentido de la evaluación}
Evaluar no es solo revisar una nota. Es reconocer qué aprendí, qué cambió en mí, qué emociones manejé, cómo me relaciono con la ciudad y la comunidad, y qué le dejaré a quien viene detrás.
\end{softbox}

\textbf{Mi reflexión liberadora en cinco dimensiones.} Completa cada frase iniciada:

\small
\begin{tabularx}{\textwidth}{>{\bfseries\RaggedRight\arraybackslash}p{3.4cm}Y>{\RaggedRight\arraybackslash}p{4.6cm}}
\rowcolor{milcVino}\color{white}Dimensión & \color{white}\bfseries Mi respuesta (frame starter) & \color{white}\bfseries Algo que descubrí\\
1. Personal & Lo que más cambió en mí fue \linead{6.5cm} & \linead{4.2cm}\\
2. Emocional & La emoción más fuerte fue \linead{2.5cm} y la regulé \linead{3cm} & \linead{4.2cm}\\
3. Ciudadana & En democracia esto importa porque \linead{6cm} & \linead{4.2cm}\\
4. Local (Cartago) & En mi barrio sería útil para \linead{6.5cm} & \linead{4.2cm}\\
5. Intergeneracional & A mi abuela/abuelo le diría \linead{6.5cm} & \linead{4.2cm}\\
\end{tabularx}
\normalsize

\begin{infoband}{milcVino}
\textbf{Para la próxima sustentación.} Vuelve a esta tabla en seis meses. Verás que las respuestas cambian — no porque lo aprendido cambie, sino porque tú cambias. La evaluación liberadora es una conversación contigo a través del tiempo.
\end{infoband}


\puente{Y para terminar, tres voces sobre los moldes y la identidad. Dussel sobre copiar vs encontrar, Epicteto sobre lo que depende de ti, Floridi sobre estándares y diversidad.}

\begin{titlebox}{milcGradeDark}
Triángulo de pensamiento · cierre filosófico
\end{titlebox}

\begin{softbox}{milcVino}{milcGris}{Dussel · lente del nosotros}
\textit{``El sur global piensa desde su propia herida y desde ahí ofrece mundo.''}

Las plantillas son productos del norte tecnológico (Silicon Valley, Europa). Usarlas tal cual te invisibiliza; adaptarlas con tu identidad local --- tus colores del Valle, tu copy en español de Colombia, tus historias reales --- es un acto pequeño de soberanía digital.

\textbf{¿Mi sitio publicado se siente como un sitio del Valle o como un clon de plantilla global? ¿Qué señales locales le agregué?}
\end{softbox}
\linead{\linewidth}\\[1mm]
\linead{\linewidth}

\begin{softbox}{milcMostaza}{milcGris}{Estoicismo · Epicteto · cuidado interior}
\textit{``Algunas cosas dependen de nosotros y otras no --- distinguir esto es la libertad.''}

Que la plantilla traiga ciertas decisiones de diseño \emph{no depende de ti}. Pero qué conservas, qué cambias, qué agregas y qué eliminas: eso sí depende. Adaptar bien es ejercer la libertad estoica sobre la propia identidad digital.

\textbf{¿Estoy ejerciendo mi libertad sobre la plantilla, o me dejé llevar por todo lo que venía? ¿Qué partes son auténticamente mías?}
\end{softbox}
\linead{\linewidth}\\[1mm]
\linead{\linewidth}

\begin{softbox}{milcTurquesa}{milcGris}{Floridi · lente de la infoesfera}
\textit{``La estandarización facilita la cooperación; la diversidad sostiene la riqueza informacional.''}

Usar plantillas es bueno: estandariza la web, hace que más gente publique. Pero si todos los sitios se ven iguales, la infoesfera se empobrece. Adaptar es contribuir a la diversidad sin abandonar el beneficio de la estandarización.

\textbf{¿Mi adaptación aporta diversidad a la infoesfera o solo replica un patrón estándar? ¿Qué de mi sitio enriquece el panorama digital común?}
\end{softbox}
\linead{\linewidth}\\[1mm]
\linead{\linewidth}

\begin{titlebox}{milcVino}
Compromiso final
\end{titlebox}

\small
\begin{tabularx}{\textwidth}{>{\RaggedRight\arraybackslash}p{3.0cm}YYY}
\rowcolor{milcVino}\color{white}\bfseries Criterio & \color{white}\bfseries Lo logré & \color{white}\bfseries En proceso & \color{white}\bfseries Necesito apoyo\\
Comprensión del tema & Explico ideas centrales con mis palabras. & Reconozco ideas pero falta precisión. & Necesito repasar conceptos base.\\
Pensamiento computacional & Apliqué los 4 pilares al diseñar mi producto. & Apliqué algunos pilares. & Necesito releer la fase 2.\\
Aplicación y producto & Mi producto cumple los criterios y se puede revisar. & Producto funcional con ajustes pendientes. & Aún en construcción.\\
Reflexión 5 dimensiones & Respondí cada dimensión con honestidad. & Respondí algunas con detalle. & Mis respuestas son muy generales.\\
Triángulo de pensamiento & Conecté las 3 voces con mi trabajo real. & Conecté al menos una. & No relaciono las voces con mi proceso.\\
\end{tabularx}
\normalsize

\textbf{Hoy entendí que:}\\[1mm]
\linead{\linewidth}\\[1mm]
\linead{\linewidth}

\textbf{Mi compromiso para la próxima sesión es:}\\[1mm]
\linead{\linewidth}\\[1mm]
\linead{\linewidth}

\begin{softbox}{milcMostaza}{milcGris}{Cita de despedida · Marco Aurelio}
\textit{``El obstáculo en el camino se vuelve el camino. Lo que estorba al camino se vuelve camino.''}\\[1mm]
Cada error de hoy, bien observado, se convierte en pieza del camino que pisarás mañana.
\end{softbox}

\small
\begin{tabularx}{\textwidth}{>{\bfseries}p{3.6cm}Y}
\rowcolor{milcGradeDark!12}Nombre completo: & \linead{10cm}\\
Fecha: & \linead{4cm} \quad Curso/grupo: \linead{4cm}\\
Mi firma: & \linead{9cm}\\
\end{tabularx}
\normalsize

\begin{infoband}{milcGradeDark}
\textbf{Para la próxima semana.} Comparte tu URL pública con 3 personas: un familiar, un compañero, un docente. Pídeles a cada uno una observación. No discutas las observaciones: anótalas. La próxima guía empezamos con esas tres voces externas.
\end{infoband}

\end{document}
